\section{The sharp bound: the parts of the proof}\label{sec:sharp}

Sections~\ref{sec:sharp} to~\ref{sec:examples} prove Theorem~\ref{thm:sharp}. This section states
the three parts of the proof and derives the theorem from them. Sections~\ref{sec:affine}
to~\ref{sec:halfproof} prove the first part, Theorem~\ref{thm:half}. Section~\ref{sec:inputs}
describes the other two parts, whose proofs are in the supplement. None of these sections is
formalised as a whole. Table~\ref{tab:sharpsteps} lists the steps that are.

\subsection{Conventions}

As in Section~\ref{sec:forest}, graphs may have loops and parallel edges. A loop adds two to the
degree of its vertex. A loop is a cycle of length one, and two parallel edges form a cycle of
length two. A multigraph is \emph{cubic} if every vertex has degree three, and \emph{subcubic} if
every vertex has degree at most three. The \emph{girth} of a multigraph is the length of a shortest
cycle; a forest has infinite girth. We write $c(G)$ for the number of cycles of $G$.

If we replace an edge by a path, the rank and the number of cycles do not change. So $\psi(r)$ is
also the largest number of cycles of a multigraph of rank $r$.

A \emph{cyclic cut} of $G$ is a set $\delta X$ of edges such that $G[X]$ and $G-X$ each contain a
cycle. Its \emph{size} is $|\delta X|$.

The two auxiliary theorems are stated for cubic multigraphs without loops.

\begin{theorem}[Small cyclic cuts; Astra, supplement~\cite{S1}]\label{thm:smallcut}
There is an absolute constant $K_4$ such that every connected cubic multigraph without loops, of
rank $r$, that has a cyclic cut of size at most four satisfies
\[
  c(G)\;\le\;2^{r-1}+K_4\,r\,2^{r/2}.
\]
\end{theorem}

\begin{theorem}[Bounded girth; Astra, supplement~\cite{S2}]\label{thm:girth}
For every integer $g_0\ge2$ there is $R(g_0)$ with the following property. Let $G$ be a connected
cubic multigraph without loops, of rank $r\ge R(g_0)$ and girth at most $g_0$, that has no cyclic
cut of size at most four. Then $c(G)<2^{r-1}$.
\end{theorem}

The graphs of Horak and Kim have cyclic cuts of size four, and they have more than $2^{r-1}$
cycles. So the second term of Theorem~\ref{thm:sharp} comes from Theorem~\ref{thm:smallcut} alone.
In the other two cases the proof gives no excess over $2^{r-1}$ beyond the excess at smaller rank.

\subsection{Cubic graphs}

\begin{lemma}\label{lem:cubic}
For every $r\ge2$ there is a connected cubic multigraph of rank $r$ with exactly $\psi(r)$ cycles.
\end{lemma}

\begin{proof}
The numbers of cycles of the graphs of rank $r$ are integers below $2^{r}$, so some graph $G_0$
attains $\psi(r)$. If $G_0$ is not connected, identify a vertex of one component with a vertex of
another. This changes neither the rank nor the set of cycles. Now delete a vertex of degree one, or
replace a vertex of degree two and its two edges by one edge, as long as this is possible. These
steps keep the rank and the number of cycles. They may produce loops and parallel edges. The result
is a connected multigraph with all degrees at least three, because $r\ge2$.

Let $v$ be a vertex of degree $d\ge4$. Replace $v$ by a path $v^{1}v^{2}\cdots v^{d-2}$, and
distribute the $d$ edge ends at $v$ along this path: two at $v^{1}$, one at each of
$v^{2},\dots,v^{d-3}$, and two at $v^{d-2}$. The new multigraph has $d-3$ more vertices and $d-3$
more edges, so its rank is $r$. A cycle through $v$ uses two edge ends at $v$. Together with the
part of the path between these two ends it is a cycle of the new multigraph. Different cycles give
different cycles, and the cycles that avoid $v$ remain. So the number of cycles does not decrease,
and it cannot increase beyond $\psi(r)$. Repeat this for every vertex of degree at least four.
\end{proof}

Entringer and Slater proved the same statement with simple cubic graphs for $r\ge3$~\cite{ES}. We
do not need it in that form.

\subsection{Proof of Theorem~\ref{thm:sharp}}

The lower bound for $\psi(r)$ is~\eqref{eq:hk}; Section~\ref{sec:examples} recalls the graphs.
For the upper bound, let $D_{*}$ and $R$ be as in Theorem~\ref{thm:half}, let $K_4$ be as in
Theorem~\ref{thm:smallcut}, and let $R(\cdot)$ be as in Theorem~\ref{thm:girth}. Choose
$R_0\ge\max\{R,\,R(D_{*}),\,9\}$ so large that $2\log(2r)\le3\log r$ for $r\ge R_0$. Put
\[
  K=\max\bigl\{K_4,\;2^{R_0/2}\bigr\}.
\]
We prove by induction on $r$ that $\psi(r)\le2^{r-1}+Kr2^{r/2}$. For $1\le r<R_0$ this holds because
$\psi(r)<2^{r}\le2^{r-1}+2^{R_0/2}2^{r/2}$. Let $r\ge R_0$, and let $G$ be a connected cubic
multigraph of rank $r$ with $\psi(r)$ cycles (Lemma~\ref{lem:cubic}). There are four cases.

\emph{Case 1: $G$ has a loop $e$.} The loop is the only cycle through $e$, and $G-e$ is connected of
rank $r-1$. So $c(G)=c(G-e)+1\le2^{r-2}+K(r-1)2^{(r-1)/2}+1\le2^{r-1}+Kr2^{r/2}$.

\emph{Case 2: $G$ has no loop, and it has a cyclic cut of size at most four.} Then
Theorem~\ref{thm:smallcut} gives the bound, because $K\ge K_4$.

\emph{Case 3: $G$ has no loop and no cyclic cut of size at most four, and its girth is less than
$D_{*}$.} Then $c(G)<2^{r-1}$ by Theorem~\ref{thm:girth} with $g_0=D_{*}$.

\emph{Case 4: $G$ has no loop and its girth is at least $D_{*}$.} Let $W$ be a shortest cycle and
let $g$ be its length. A cubic multigraph of girth $g$ has at least $2^{g/2}$ vertices, and $G$ has
$2r-2$ vertices. So $g\le2\log(2r)\le3\log r$. The set $P=V(W)$ is connected, and $G-P$ has no cycle
of length less than $D_{*}$. So $q(G,P)\le\half$ by Theorem~\ref{thm:half}. Let $e$ be an edge of
$W$. By Lemma~\ref{lem:edge} and the induction hypothesis for the connected multigraph $G-e$,
\[
  c(G)\;\le\;c(G-e)+2^{r-2}\;\le\;2^{r-1}+K(r-1)2^{(r-1)/2}\;\le\;2^{r-1}+Kr2^{r/2}.
\]
This completes the induction. A graph of rank $r$ that is not connected has at most $\psi(r)$
cycles by definition. \qed

\begin{remark}
In Case~3 the girth is at least five when $r\ge9$, although the proof does not need this. Let $C$
be a cycle of length at most four in a cubic multigraph without loops, and put $X=V(C)$. At most
eight edges have an end in $X$, and at most four of them leave $X$. So $G-X$ has rank at least
$r-8$. If $r\ge9$, then $\delta X$ is a cyclic cut of size at most four.
\end{remark}

\begin{remark}
The multigraph $G-e$ in Case~4 is not cubic. This is why the induction is on $\psi(r)$ and not on a
class of cubic graphs.
\end{remark}
